% GENERATED FILE. Edit canonical lessons, then run npm run generate:books.
% canonical-source-hash: fnv1a64:9456f0340081776c
% canonical-lessons: KA-C183-banale, KA-C183-kukkar, KA-C183-battalu, KA-C183-sautu, KA-C183-baket

\chapter{Frying pan, Pressure cooker, Bowl, Ladle, Bucket}
\label{ch:ka-frying-pan-pressure-cooker-bowl-ladle-bucket}
\hlchaptermodality{183}

\begin{chapteropening}
\textbf{By the end of this chapter:} \emph{I can say a frying pan, a pressure cooker, a bowl, a ladle and a bucket.}

Complete the last lesson of chapter 183: I can say a frying pan, a pressure cooker, a bowl, a ladle and a bucket.
\end{chapteropening}

\section[\texorpdfstring{\kn{ಬಾಣಲೆ} (bāṇale)}{ಬಾಣಲೆ (bāṇale)}]{\kn{ಬಾಣಲೆ} (bāṇale) — a frying pan}

\label{lesson:KA-C183-banale}



\begin{quote}
Before the new one: say the Kannada for ripe, then the Kannada for one's own.
\end{quote}

\subsection*{You'll want to know: \kn{ಬಾಣಲೆ}}
\textbf{\kn{ಬಾಣಲೆ}} — \emph{bāṇale} — \textquotedblleft{}a frying pan\textquotedblright{}.

\begin{grammarlens}[title={What you've built}]
One more word for this chapter.
\end{grammarlens}

\subsection*{Your turn}
\begin{itemize}
\raggedright
  \item \emph{Say it:} \emph{bāṇale}
  \item \emph{Say it:} \emph{bāṇale}, once more
  \item \emph{Say it:} say \emph{bāṇale}
  \item \emph{Recall:} say the Kannada for ripe, then the Kannada for one's own, then say \emph{bāṇale} again
\end{itemize}

\begin{tcolorbox}[breakable,colback=teal!4,colframe=teal!35!black,title={Before you move on}]
What does \kn{ಬಾಣಲೆ} mean? (\textquotedblleft{}A frying pan\textquotedblright{}.) Say it once more.
\end{tcolorbox}
\section[\texorpdfstring{\kn{ಕುಕ್ಕರ್} (kukkar)}{ಕುಕ್ಕರ್ (kukkar)}]{\kn{ಕುಕ್ಕರ್} (kukkar) — a pressure cooker}

\label{lesson:KA-C183-kukkar}



\begin{quote}
Before the new one: say the Kannada for one's own, then the Kannada for a frying pan.
\end{quote}

\subsection*{You'll want to know: \kn{ಕುಕ್ಕರ್}}
\textbf{\kn{ಕುಕ್ಕರ್}} — \emph{kukkar} — \textquotedblleft{}a pressure cooker\textquotedblright{}.

\begin{grammarlens}[title={What you've built}]
One more word for this chapter.
\end{grammarlens}

\subsection*{Your turn}
\begin{itemize}
\raggedright
  \item \emph{Say it:} \emph{kukkar}
  \item \emph{Say it:} \emph{kukkar}, once more
  \item \emph{Say it:} say \emph{kukkar}
  \item \emph{Recall:} say the Kannada for one's own, then the Kannada for a frying pan, then say \emph{kukkar} again
\end{itemize}

\begin{tcolorbox}[breakable,colback=teal!4,colframe=teal!35!black,title={Before you move on}]
What does \kn{ಕುಕ್ಕರ್} mean? (\textquotedblleft{}A pressure cooker\textquotedblright{}.) Say it once more.
\end{tcolorbox}
\section[\texorpdfstring{\kn{ಬಟ್ಟಲು} (baṭṭalu)}{ಬಟ್ಟಲು (baṭṭalu)}]{\kn{ಬಟ್ಟಲು} (baṭṭalu) — a bowl}

\label{lesson:KA-C183-battalu}



\begin{quote}
Before the new one: say the Kannada for a frying pan, then the Kannada for a pressure cooker.
\end{quote}

\subsection*{You'll want to know: \kn{ಬಟ್ಟಲು}}
\textbf{\kn{ಬಟ್ಟಲು}} — \emph{baṭṭalu} — \textquotedblleft{}a bowl\textquotedblright{}.

\begin{grammarlens}[title={What you've built}]
One more word for this chapter.
\end{grammarlens}

\subsection*{Your turn}
\begin{itemize}
\raggedright
  \item \emph{Say it:} \emph{baṭṭalu}
  \item \emph{Say it:} \emph{baṭṭalu}, once more
  \item \emph{Say it:} say \emph{baṭṭalu}
  \item \emph{Recall:} say the Kannada for a frying pan, then the Kannada for a pressure cooker, then say \emph{baṭṭalu} again
\end{itemize}

\begin{tcolorbox}[breakable,colback=teal!4,colframe=teal!35!black,title={Before you move on}]
What does \kn{ಬಟ್ಟಲು} mean? (\textquotedblleft{}A bowl\textquotedblright{}.) Say it once more.
\end{tcolorbox}
\section[\texorpdfstring{\kn{ಸೌಟು} (sauṭu)}{ಸೌಟು (sauṭu)}]{\kn{ಸೌಟು} (sauṭu) — a ladle}

\label{lesson:KA-C183-sautu}



\begin{quote}
Before the new one: say the Kannada for a pressure cooker, then the Kannada for a bowl.
\end{quote}

\subsection*{You'll want to know: \kn{ಸೌಟು}}
\textbf{\kn{ಸೌಟು}} — \emph{sauṭu} — \textquotedblleft{}a ladle\textquotedblright{}.

\begin{grammarlens}[title={What you've built}]
One more word for this chapter.
\end{grammarlens}

\subsection*{Your turn}
\begin{itemize}
\raggedright
  \item \emph{Say it:} \emph{sauṭu}
  \item \emph{Say it:} \emph{sauṭu}, once more
  \item \emph{Say it:} say \emph{sauṭu}
  \item \emph{Recall:} say the Kannada for a pressure cooker, then the Kannada for a bowl, then say \emph{sauṭu} again
\end{itemize}

\begin{tcolorbox}[breakable,colback=teal!4,colframe=teal!35!black,title={Before you move on}]
What does \kn{ಸೌಟು} mean? (\textquotedblleft{}A ladle\textquotedblright{}.) Say it once more.
\end{tcolorbox}
\section[\texorpdfstring{\kn{ಬಕೆಟ್} (bakeṭ)}{ಬಕೆಟ್ (bakeṭ)}]{\kn{ಬಕೆಟ್} (bakeṭ) — a bucket}

\label{lesson:KA-C183-baket}



\begin{quote}
Before the new one: say the Kannada for a bowl, then the Kannada for a ladle.
\end{quote}

\subsection*{You'll want to know: \kn{ಬಕೆಟ್}}
\textbf{\kn{ಬಕೆಟ್}} — \emph{bakeṭ} — \textquotedblleft{}a bucket\textquotedblright{}.

\begin{grammarlens}[title={What you've built}]
One more word for this chapter.
\end{grammarlens}

\subsection*{Your turn}
\begin{itemize}
\raggedright
  \item \emph{Say it:} \emph{bakeṭ}
  \item \emph{Say it:} \emph{bakeṭ}, once more
  \item \emph{Say it:} say \emph{bakeṭ}
  \item \emph{Recall:} say the Kannada for a bowl, then the Kannada for a ladle, then say \emph{bakeṭ} again
\end{itemize}

\begin{tcolorbox}[breakable,colback=teal!4,colframe=teal!35!black,title={Before you move on}]
What does \kn{ಬಕೆಟ್} mean? (\textquotedblleft{}A bucket\textquotedblright{}.) Say it once more.
\end{tcolorbox}
